\appendix

\chapter{Constants and their provenance}
\label{app:constants}

Every number used in this document, with where it came from. The status
column distinguishes a measurement from a calculation from an assumption,
because the three carry different weight and the distinction is lost once
they share a table in a slide.

\begin{longtable}{@{}p{0.38\textwidth} p{0.17\textwidth} p{0.38\textwidth}@{}}
\caption[Constants used]{Constants and derived quantities used in this
document.}
\label{tab:constants} \\
\toprule
\textbf{Quantity} & \textbf{Value} & \textbf{Status and source} \\
\midrule
\endfirsthead
\toprule
\textbf{Quantity} & \textbf{Value} & \textbf{Status and source} \\
\midrule
\endhead
\multicolumn{3}{@{}l}{\textit{The Sun}} \\
Total solar irradiance, AM0 & 1361~W/m\textsuperscript{2} & IAU 2015 nominal
  \cite{prsa2016iau} \\
Solar illuminance in orbit & 132\,800~lx & Calculated, AM0 spectrum against
  $V(\lambda)$ \cite{astm_e490_data, cie_vlambda} \\
Solar effective temperature & 5772~K & IAU 2015 nominal \cite{prsa2016iau} \\
Solar angular diameter & 0.533\textdegree & IAU radius and distance
  \cite{prsa2016iau} \\
Penumbra cast by the Sun & 9.3~mm per m & Calculated from the above,
  Eq. \ref{eq:penumbra} \\
\addlinespace
\multicolumn{3}{@{}l}{\textit{Sensor and lens, as fitted}} \\
Sensor & Sony IMX477 & Hardware \cite{sony_imx477_datasheet} \\
Pixel pitch & 1.55~\textmu m & Datasheet \cite{sony_imx477_datasheet} \\
Full sensor resolution & $4056 \times 3040$ & Datasheet \\
Focal length & 8~mm & Lens \cite{arducam_c1508zm04} \\
Aperture range & f/1.6 to f/16 & Lens \cite{arducam_c1508zm04} \\
Horizontal field of view & 42.9\textdegree & Calculated from active area and
  focal length. \textbf{Not} the 58.4\textdegree{} on the lens datasheet,
  which is for a 2/3'' format \\
Stereo baseline & 100~mm & Integrated support \\
Fitted $f_x$, best archived session & $2651.79 \pm 22.42$~px & Measured, at
  $1920 \times 1440$ \cite{linku_calibration_method} \\
\addlinespace
\multicolumn{3}{@{}l}{\textit{Scene}} \\
Nominal working distance & 3.5~m & From the Blender replica, not a measured
  room dimension \\
Tether diameter, physical proxy & 1.0~mm & Laboratory item \\
Tether diameter, analytical worst case & 0.8~mm & \textbf{Assumption} stated
  as such in \cite{linku_d5_camera_report}. Do not confuse with the proxy \\
Tether length, mission design figure & up to 100~m & Mission design, not
  measured \\
\addlinespace
\multicolumn{3}{@{}l}{\textit{Derived optics at 3.5~m}} \\
Projected tether width & 1.47~px & Calculated \cite{linku_dof_visualizer} \\
Depth of field at f/1.6 & 3.08 to 4.05~m & Calculated, 1~px circle of
  confusion \cite{linku_dof_visualizer} \\
Hyperfocal distance at f/1.6 & 25.8~m & Calculated \\
Airy disc at f/1.6 & 2.15~\textmu m, 1.39~px & Calculated \\
\addlinespace
\multicolumn{3}{@{}l}{\textit{Motion and exposure}} \\
Reference angular rate & 1.0\textdegree/s & Corresponds to the software's
  current 15~frame/s configuration, \textbf{not} a sensor limit
  \cite{linku_rollingshutter} \\
Sensor-limited angular rate & 2.76\textdegree/s & 12-bit full-resolution
  readout. Itself a \textbf{bound}: the exact row-scan time needs the
  IMX477 Software Reference Manual, which is not public \\
Shear budget & 3~px & \textbf{Assumption}, chosen because the tether is 1 to
  3~px wide \\
Blur budget, this document & 1~px & Project choice, deliberately stricter
  than the 3~px above \\
Maximum exposure at 1\textdegree/s & 11.1~ms & Calculated \\
\addlinespace
\multicolumn{3}{@{}l}{\textit{Illumination requirements}} \\
Minimum illuminance, fitted configuration & 2320~lx & Derived from the
  datasheet; a \textbf{floor} \cite{sony_imx477_datasheet,
  linku_darkroom_illumination} \\
Minimum illuminance, Hawk-eye & 3620~lx & \textbf{Assumption}: the IMX477
  saturation rescaled by aperture, because Arducam does not identify the
  sensor die \\
Penumbra limit & 35~mm at 1~m & Project decision, near the ESA
  1.9\textdegree{} reference \cite{esa_lss} \\
Uniformity limit & 5\% & IEC 60904-9 class B \cite{iec60904_9_2020} \\
Stability limit & 2\% & IEC 60904-9 class B short term \\
Colour temperature & 5600 to 6000~K & Project decision, from laboratory
  practice \cite{dlr2025eposportfolio} \\
\addlinespace
\multicolumn{3}{@{}l}{\textit{Calibration target}} \\
Board & ChArUco $5 \times 7$ & \texttt{DICT\_6X6\_250}
  \cite{garrido2014aruco} \\
Square and marker size & 37~mm, 29.6~mm & Generated by the repository tool \\
Tilt bins & 0--10, 10--30, 30--45, 45--70\textdegree & Project split,
  following the $\pm70$\textdegree{} range of \cite{rojtberg2018pose} \\
Minimum views per bin & 3 & Project decision \\
\bottomrule
\end{longtable}

\chapter{Sources, with their provenance level}
\label{app:sources}

Each cited source is archived in the \texttt{Sources/} folder of this
deliverable. Where a source could not be obtained in full, or is secondary,
that is stated here rather than buried in a footnote, because a reader
checking a number is entitled to know how firm the ground under it is.

\begin{longtable}{@{}p{0.30\textwidth} p{0.14\textwidth} p{0.49\textwidth}@{}}
\caption[Source provenance]{Provenance of the cited sources.}
\label{tab:sources} \\
\toprule
\textbf{Source} & \textbf{Level} & \textbf{Note} \\
\midrule
\endfirsthead
\toprule
\textbf{Source} & \textbf{Level} & \textbf{Note} \\
\midrule
\endhead
IEC 60904-9:2020 \cite{iec60904_9_2020} & \textbf{Partial} & Only the free
  official preview was read, clauses 1 to 5.2, which does include all three
  class tables. The measurement procedures after 5.2 remain behind the
  paywall. Every class threshold adopted rests on this partial read \\
ASTM E490 AM0 spectrum \cite{astm_e490_data} & Primary, mirrored & The data
  file is taken from a public mirror; the original NREL path is dead \\
CIE $V(\lambda)$ \cite{cie_vlambda} & Primary & \\
IAU 2015 nominal solar values \cite{prsa2016iau} & Primary & \\
Sony IMX477 datasheet \cite{sony_imx477_datasheet} & Primary & Sections 10-1
  and 11-3/11-4 give the sensitivity and saturation figures used. The
  Software Reference Manual, which holds the exact readout timing, is not
  public \\
Arducam C1508ZM04 \cite{arducam_c1508zm04} & Vendor & Its quoted
  58.4\textdegree{} horizontal field of view applies to a 2/3'' format and
  must not be used with this sensor \\
Raspberry Pi camera documentation \cite{raspberrypi_camera_docs} & Primary &
  Source of the sensor mode behaviour, including that some modes crop as
  well as bin \\
ARRI EB 12/18 manual \cite{arri_eb1218_manual} & Primary & The statement
  that the 50/60~Hz mode is not flicker free is the manufacturer's own \\
EBU Tech 3355, TLCI \cite{ebu_tech3355} & Primary & \\
Sheinin et al., rolling shutter on the grid \cite{sheinin2018rollingac} &
  Primary & \\
ESA Large Space Simulator \cite{esa_lss} & Primary, web & Beam
  specifications from the agency's own page \\
DLR EPOS portfolio \cite{dlr2025eposportfolio} & Primary & Operator
  document. Note that a widely cited survey gives the fixture designation
  incorrectly as ARRI Max 18/12; the operator documents say 12/18 \\
Benninghoff et al., EPOS 2.0 \cite{benninghoff2017epos} & Primary & \\
Park et al., TRON \cite{park2021tron} & Primary & Operator paper. A web
  summary describes the lamp as a xenon short arc; the operators state a
  metal halide arc \\
Pauly et al., lessons from a space lab \cite{pauly2022spacelab} & Preprint &
  The arXiv v2 text is the version read, which differs from the later
  journal version \\
Zhang, flexible camera calibration \cite{zhang2000flexible} & Primary & \\
Rojtberg and Kuijper, pose selection \cite{rojtberg2018pose} & Primary &
  Source of the tilt range the bins cover \\
Garrido-Jurado et al., ArUco \cite{garrido2014aruco} & Primary & \\
OpenCV calib3d \cite{opencv_calib3d} & Primary & \\
LINKU companion deliverables
  \cite{linku_d5_camera_report, linku_darkroom_illumination,
  linku_calibration_method, linku_rollingshutter,
  linku_blender_validation, linku_workplan} & Internal & Two of these are
  working drafts whose results chapters are pending; the camera system
  report contains no reconstruction accuracy figure \\
\texttt{camera\_dof\_visualizer.py} \cite{linku_dof_visualizer} & Internal
  tool & Thin-lens model. The depth-of-field figures quoted in this document
  were reproduced independently from the thin-lens equations as a check \\
\bottomrule
\end{longtable}
